% Don't like 10pt? Try 11pt or 12pt
\documentclass[10pt]{article}

% The automated optical recognition software used to digitize resume
% information works best with fonts that do not have serifs. This
% command uses a sans serif font throughout. Uncomment both lines (or at
% least the second) to restore a Roman font (i.e., a font with serifs).
%\usepackage{times}
%\renewcommand{\familydefault}{\sfdefault}

% This is a helpful package that puts math inside length specifications
\usepackage{calc}
\usepackage{comment}
\usepackage{color}
% Simpler bibsection for CV sections
% (thanks to natbib for inspiration)
\makeatletter
\newlength{\bibhang}
\setlength{\bibhang}{1em} %1em}
\newlength{\bibsep}
 {\@listi \global\bibsep\itemsep \global\advance\bibsep by\parsep}
\newenvironment{bibsection}%
        {\begin{enumerate}{}{%
%        {\begin{list}{}{%
       \setlength{\leftmargin}{\bibhang}%
       \setlength{\itemindent}{-\leftmargin}%
       \setlength{\itemsep}{\bibsep}%
       \setlength{\parsep}{\z@}%
        \setlength{\partopsep}{0pt}%
        \setlength{\topsep}{0pt}}}
        {\end{enumerate}\vspace{-.6\baselineskip}}
%        {\end{list}\vspace{-.6\baselineskip}}
\makeatother

% Layout: Puts the section titles on left side of page
\reversemarginpar

%
%         PAPER SIZE, PAGE NUMBER, AND DOCUMENT LAYOUT NOTES:
%
% The next \usepackage line changes the layout for CV style section
% headings as marginal notes. It also sets up the paper size as either
% letter or A4. By default, letter was used. If A4 paper is desired,
% comment out the letterpaper lines and uncomment the a4paper lines.
%
% As you can see, the margin widths and section title widths can be
% easily adjusted.
%
% ALSO: Notice that the includefoot option can be commented OUT in order
% to put the PAGE NUMBER *IN* the bottom margin. This will make the
% effective text area larger.
%
% IF YOU WISH TO REMOVE THE ``of LASTPAGE'' next to each page number,
% see the note about the +LP and -LP lines below. Comment out the +LP
% and uncomment the -LP.
%
% IF YOU WISH TO REMOVE PAGE NUMBERS, be sure that the includefoot line
% is uncommented and ALSO uncomment the \pagestyle{empty} a few lines
% below.
%

%% Use these lines for letter-sized paper
\usepackage[paper=letterpaper,
            %includefoot, % Uncomment to put page number above margin
            marginparwidth=1.2in,     % Length of section titles
            marginparsep=.05in,       % Space between titles and text
            margin=0.75in,               % 1 inch margins
            includemp]{geometry}

%% Use these lines for A4-sized paper
%\usepackage[paper=a4paper,
%            %includefoot, % Uncomment to put page number above margin
%            marginparwidth=30.5mm,    % Length of section titles
%            marginparsep=1.5mm,       % Space between titles and text
%            margin=25mm,              % 25mm margins
%            includemp]{geometry}

%% More layout: Get rid of indenting throughout entire document
\setlength{\parindent}{0in}

\usepackage[shortlabels]{enumitem}

%% Reference the last page in the page number
%
% NOTE: comment the +LP line and uncomment the -LP line to have page
%       numbers without the ``of ##'' last page reference)
%
% NOTE: uncomment the \pagestyle{empty} line to get rid of all page
%       numbers (make sure includefoot is commented out above)
%
\usepackage{fancyhdr,lastpage}
\pagestyle{fancy}
%\pagestyle{empty}      % Uncomment this to get rid of page numbers
\fancyhf{}\renewcommand{\headrulewidth}{0pt}
\fancyfootoffset{\marginparsep+\marginparwidth}
\newlength{\footpageshift}
\setlength{\footpageshift}
          {0.5\textwidth+0.5\marginparsep+0.5\marginparwidth-2in}
\lfoot{\hspace{\footpageshift}%
       \parbox{4in}{\, \hfill %
                    \arabic{page} of \protect\pageref*{LastPage} % +LP
%                    \arabic{page}                               % -LP
                    \hfill \,}}

% Finally, give us PDF bookmarks
\usepackage{color,hyperref}
\definecolor{darkblue}{rgb}{0.0,0.0,0.3}
\hypersetup{colorlinks,breaklinks,
            linkcolor=darkblue,urlcolor=darkblue,
            anchorcolor=darkblue,citecolor=darkblue}

%%%%%%%%%%%%%%%%%%%%%%%% End Document Setup %%%%%%%%%%%%%%%%%%%%%%%%%%%%

\usepackage{CJKutf8}

%%%%%%%%%%%%%%%%%%%%%%%%%%% Helper Commands %%%%%%%%%%%%%%%%%%%%%%%%%%%%

% The title (name) with a horizontal rule under it
% (optional argument typesets an object right-justified across from name
%  as well)
%
% Usage: \makeheading{name}
%        OR
%        \makeheading[right_object]{name}
%
% Place at top of document. It should be the first thing.
% If ``right_object'' is provided in the square-braced optional
% argument, it will be right justified on the same line as ``name'' at
% the top of the CV. For example:
%
%       \makeheading[\emph{Curriculum vitae}]{Your Name}
%
% will put an emphasized ``Curriculum vitae'' at the top of the document
% as a title. Likewise, a picture could be included:
%
%   \makeheading[\includegraphics[height=1.5in]{my_picutre}]{Your Name}
%
% the picture will be flush right across from the name.
\newcommand{\makeheading}[2][]%
        {\hspace*{-\marginparsep minus \marginparwidth}%
         \begin{minipage}[t]{\textwidth+\marginparwidth+\marginparsep}%
             {\large \bfseries #2 \hfill #1}\\[-0.15\baselineskip]%
                 \rule{\columnwidth}{1pt}%
         \end{minipage}}

% The section headings
%
% Usage: \section{section name}
\renewcommand{\section}[1]{\pagebreak[3]%
    \hyphenpenalty=10000%
    \vspace{1.3\baselineskip}%
    \phantomsection\addcontentsline{toc}{section}{#1}%
    \noindent\llap{\scshape\smash{\parbox[t]{\marginparwidth}{\raggedright #1}}}%
    \vspace{-\baselineskip}\par}

% An itemize-style list with lots of space between items
\newenvironment{outerlist}[1][\enskip\textbullet]%
        {\begin{itemize}[#1,leftmargin=*]}{\end{itemize}%
         \vspace{-.6\baselineskip}}

% An environment IDENTICAL to outerlist that has better pre-list spacing
% when used as the first thing in a \section
\newenvironment{lonelist}[1][\enskip\textbullet]%
        {\begin{list}{#1}{%
        \setlength{\partopsep}{0pt}%
        \setlength{\topsep}{0pt}}}
        {\end{list}\vspace{-.6\baselineskip}}

% An itemize-style list with little space between items
\newenvironment{innerlist}[1][\enskip\textbullet]%
        {\begin{itemize}[#1,leftmargin=*,parsep=0pt,itemsep=0pt,topsep=0pt,partopsep=0pt]}
        {\end{itemize}}

% An environment IDENTICAL to innerlist that has better pre-list spacing
% when used as the first thing in a \section
\newenvironment{loneinnerlist}[1][\enskip\textbullet]%
        {\begin{itemize}[#1,leftmargin=*,parsep=0pt,itemsep=0pt,topsep=0pt,partopsep=0pt]}
        {\end{itemize}\vspace{-.6\baselineskip}}

% To add some paragraph space between lines.
% This also tells LaTeX to preferably break a page on one of these gaps
% if there is a needed pagebreak nearby.
\newcommand{\blankline}{\quad\pagebreak[3]}
\newcommand{\halfblankline}{\quad\vspace{-0.5\baselineskip}\pagebreak[3]}

% Uses hyperref to link DOI
\newcommand\doilink[1]{\href{http://dx.doi.org/#1}{#1}}
\newcommand\doi[1]{doi:\doilink{#1}}

% For \url{SOME_URL}, links SOME_URL to the url SOME_URL
\providecommand*\url[1]{\href{#1}{#1}}
% Same as above, but pretty-prints SOME_URL in teletype fixed-width font
\renewcommand*\url[1]{\href{#1}{\texttt{#1}}}

% For \email{ADDRESS}, links ADDRESS to the url mailto:ADDRESS
\providecommand*\email[1]{\href{mailto:#1}{#1}}
% Same as above, but pretty-prints ADDRESS in teletype fixed-width font
%\renewcommand*\email[1]{\href{mailto:#1}{\texttt{#1}}}

%\providecommand\BibTeX{{\rm B\kern-.05em{\sc i\kern-.025em b}\kern-.08em
%    T\kern-.1667em\lower.7ex\hbox{E}\kern-.125emX}}
%\providecommand\BibTeX{{\rm B\kern-.05em{\sc i\kern-.025em b}\kern-.08em
%    \TeX}}
\providecommand\BibTeX{{B\kern-.05em{\sc i\kern-.025em b}\kern-.08em
    \TeX}}
\providecommand\Matlab{\textsc{Matlab}}

%%%%%%%%%%%%%%%%%%%%%%%% End Helper Commands %%%%%%%%%%%%%%%%%%%%%%%%%%%

%%%%%%%%%%%%%%%%%%%%%%%%% Begin CV Document %%%%%%%%%%%%%%%%%%%%%%%%%%%%

\begin{document}
\makeheading{Pengzhi Gao}

\section{Contact Information}
% NOTE: Mind where the & separators and \\ breaks are in the following
%       table.
%
% ALSO: \rcollength is the width of the right column of the table
%       (adjust it to your liking; default is 1.85in).
%
\newlength{\rcollength}\setlength{\rcollength}{2.5in}
\begin{tabular}[t]{@{}p{\textwidth-\rcollength}p{\rcollength}}
Xiaomi Campus & \textit{Mobile: } (+86) 15810512592 \\
No. 33 Xierqi Middle Road, & \textit{Email: } gpengzhi@gmail.com\\
Haidian District, Beijing, China 100085  & \textit{Homepage: } https://gpengzhi.github.io/\\
\end{tabular}

\section{Education}

{\textbf{Rensselaer Polytechnic Institute}}, Troy, NY, USA
\begin{outerlist}
\item[] Ph.D., {Electrical Engineering}, August 2013 - December 2017 
% \begin{innerlist}
% \item Advisor: {Professor Meng Wang}
% \item Thesis: High-dimensional Data Analysis by Exploiting Low-dimensional Models with Applications in Synchrophasor Data Analysis in Power Systems
% \end{innerlist}
\end{outerlist}

\vspace{.1in}
{\textbf{University of Pennsylvania}}, Philadelphia, PA, USA
\begin{outerlist}
\item[] M.S., {Electrical Engineering}, September 2011 - May 2013 
% \begin{innerlist}
% \item GPA: 3.74/4
% \end{innerlist}
\end{outerlist}

\vspace{.1in}
{\textbf{Xidian University}}, Xi'an, China
\begin{outerlist}
\item[] B.S. (with honors), {Electronic and Information Engineering}, August 2007 - July 2011
% \begin{innerlist}
% \item GPA: 91.4/100 \textbf{ }\textbf{ } Class Rank: 1st in 114
% \end{innerlist}
\end{outerlist}

%\section{Research Interests}

%My research interests lie in the intersection of the fields of signal processing, high-dimensional statistics, and natural language processing, with a particular emphasis on developing optimization methods and training strategies with applications in smart grids, machine translation, and multilingual large language models.

% I am particularly interested in developing low-dimensional models and optimization methods with applications in image processing and power system monitoring.

\section{Work Experience}

\textbf{Senior Algorithm Engineer \& Tech Lead Manager} \hfill {June 2024 to Present}
\begin{innerlist}
\item[] Large Language Model Team,\\
\textbf{Xiaomi, Inc.}, Beijing, China \\
Supervisor: Wei Liu and Jian Luan
\item Led a cross-functional UI Agent \& Agentic RL team of fifteen researchers, engineers, and interns to build native GUI Agent models that deliver one-step direct action capabilities in XiaoAi, enabling the assistant to autonomously complete multi-step GUI tasks from user instructions.
\item Developed many-to-many multilingual machine translation models, GemmaX (NAACL 2025), which became the top trending model in Hugging Face's machine translation category, gaining over two hundred thousand downloads, and supporting real-time subtitle translation systems in XiaoAi Translate.
\end{innerlist}

\

\textbf{Staff Research \& Development Engineer} \hfill {September 2020 to June 2024}
\begin{innerlist}
\item[] Natural Language Processing Department,\\
\textbf{Baidu, Inc.}, Beijing, China \\
Supervisor: Dr. Zhongjun He
% \item Investigated and proposed consistency-based training strategy to boost the performance of end-to-end speech-to-text translation (NAACL 2024). Investigated and proposed cross-lingual consistency-based training strategy to boost many-to-many multilingual machine translation with large language models (arXiv 2024).
% \item Developed multilingual machine translation system based on large language models such as LLaMA-2, Qwen1.5, and Ernie models. Investigated cross-lingual consistency-based training strategy for boosting many-to-many multilingual translation with LLMs (arXiv 2024).
\item Developed a unified multilingual sentence representation model, MuSR, that supports more than 220 languages (EMNLP 2023 Industry Track), which is used as the foundation model for the similarity-based bitext mining system in Baidu Translate.
\item Proposed a simple but effective training strategy to boost the zero-shot performance of multilingual neural machine translation (ACL 2023 Findings). Developed data processing and model training pipelines and improved the multilingual machine translation system that supports more than 200 languages in Baidu Translate.
\item Designed Bi-SimCut, a simple but effective training strategy that boosts the performance of neural machine translation (NAACL 2022). Developed a training pipeline based on Bi-SimCut and improved the machine translation system in Baidu Translate.
% \item Proposed a novel method, MixDiversity, to generate different translations with high faithfulness and diversity (EMNLP 2021 Findings). Developed data augmentation method based on MixDiversity and improved the machine translation system in Baidu Translate.
\end{innerlist}

\

\textbf{Data Scientist / Machine Learning Engineer} \hfill {February 2018 to April 2020}
\begin{innerlist}
\item[] Core Machine Learning Team,\\
\textbf{Petuum, Inc.}, Pittsburgh, PA, USA \\
Supervisor: Dr. Tong Wen and Dr. Zhiting Hu
\item Designed and implemented a machine learning library for Petuum AI Builder Platform.
\item Designed and developed Texar-PyTorch (gaining over 745 stars on GitHub), an open-source machine learning and text generation toolkit based on PyTorch. Maintained and contributed to Texar (gaining over 2390 stars on GitHub), an open-source machine learning and text generation toolkit based on TensorFlow.
\item Designed and developed Forte (gaining over 250 stars on GitHub), a toolkit for building natural language processing pipelines, featuring cross-task interaction, adaptable data-model interfaces and many more (EMNLP 2020 System Demonstrations).
\end{innerlist}

\

\textbf{Research Intern} \hfill {December 2010 to May 2011}
\begin{innerlist}
\item[] Internet Media Group,\\
\textbf{Microsoft Research Asia}, Beijing, China \\
Supervisor: Dr. Feng Wu and Dr. Chong Luo
\item Analyzed the data collected from 54 sensors deployed in Intel Berkeley Research Lab to exploit the temporal correlations in sensor readings. Developed a joint source network coding scheme for approximate data gathering in the wireless sensor network.
\end{innerlist}

\section{Skills}

\begin{innerlist}
\item \textbf{Research}: Generative AI, LLM/VLM-based Agents, Machine Learning, Natural Language Processing, Signal Processing, High-dimensional Statistics
\end{innerlist}
\begin{innerlist}
\item \textbf{Programming}: Python $>$ C/C++ $=$ Matlab $>$ Java $=$ R
\end{innerlist}
\begin{innerlist}
\item \textbf{Software}: PyTorch, TensorFlow, Keras, Scikit-learn, Linux, MacOS, Git
\end{innerlist}


\section{Honors and Awards}

\begin{innerlist}
\item North America Finalist of IBM Watson Build Challenge\hfill 2017
\item Paper selected as runner-up for the Best Paper in Electric Energy Systems\\ Track at Hawaii International Conference on System Sciences \hfill 2015
\item Founders Award of Excellence (top 1\%)\hfill 2015
\item Paper selected as one of the Best Conference Papers on Power System Analysis\\ and Modeling at
IEEE Power \& Energy Society General Meeting \hfill 2014
\item Excellent Graduate of Xidian University (top 1\%)\hfill 2011
\item National Scholarship (top 0.2\% nationwide)\hfill 2010
\item First prize of the College Academic and Technological Scholarship (top 2\%)\hfill 2008 - 2010
\item Excellent Student Award (top 1\%)\hfill 2008
\end{innerlist}


\section{Preprints \& Technical Reports}
$^*$ denotes equal contribution.
\vspace{-.1275in}
\begin{bibsection}

\item \href{}{G. Liu, J. Ye, {\bf P. Gao}, W. Liu, J. Luan, Y. Liu, and Y. Li. ``SimuWoB: Simulating Real-World Mobile Apps for Fast and Faithful GUI Agent Benchmarking.'' \emph{arXiv: 2605.25160}, 2026.}

\item \href{}{W. Xu, K. Huang, Y. Feng, J. Li, Y. Chen, Y. Liu, Z. Jiang, H. Qu, {\bf P. Gao}, W. Liu, J. Luan, X. Hu, and B. An. ``How Mobile World Model Guides GUI Agents?'' \emph{arXiv: 2605.10347}, 2026.}

\item \href{}{Y. Chen$^*$, {\bf P. Gao}$^*$, Q. Wu, K. Huang, Y. Liu, H. Qu, K. Shao, and J. Luan. ``Digital Agents Meet World Models: A Survey.'' May, 2026.}

\item \href{}{Y. Shang, {\bf P. Gao}, Y. Yang, J. Ma, W. Liu, J. Luan, and J. Su. ``ExPosST: Explicit Positioning with Adaptive Masking for LLM-Based Simultaneous Machine Translation.'' \emph{arXiv: 2603.14903}, 2026.}

\item \href{}{Y. Shang$^*$, {\bf P. Gao}$^*$, W. Liu, J. Luan, and J. Su. ``Scaling Model and Data for Multilingual Machine Translation with Open Large Language Models.'' \emph{arXiv: 2602.11961}, 2026.}

\item \href{}{T. Xiong, X. Hu, Y. Chen, Y. Liu, C. Wu, {\bf P. Gao}, W. Liu, J. Luan, and S. Zhang. ``GUI-PRA: Process Reward Agent for GUI Tasks.'' \emph{arXiv: 2509.23263}, 2025.}

\item \href{}{L. Gao, L. Zhang, S. Wang, {\bf P. Gao}, W. Liu, J. Luan, S. Wang, Y. Li, M. Xu. ``MobileViews: A Million-scale and Diverse Mobile GUI Dataset.'' \emph{arXiv: 2409.14337}, 2024.}

\item \href{}{{\bf P. Gao}, Z. He, H. Wu, and H. Wang. ``Towards Boosting Many-to-Many Multilingual Machine Translation with Large Language Models.'' \emph{arXiv: 2401.05861}, 2024.}

\item \href{}{R. Wang$^*$, {\bf P. Gao}$^*$, and M. Wang. ``Robust Matrix Completion by Exploiting Dynamic Low-dimensional Structures.'' \emph{DOI: 10.21203/rs.3.rs-420556/v1}, 2021.}

\item \href{}{Z. Hu, {\bf P. Gao}, A. Bukkittu, and Z. Hu. ``Introducing Texar-PyTorch: An ML Library integrating the best of TensorFlow into PyTorch.'' October, 2019.}
	
\end{bibsection}


\section{Journal Publications}
\vspace{-.1275in}
\begin{bibsection}

\item \href{}{{\bf P. Gao}$^*$, R. Wang$^*$, M. Wang, and J. H. Chow. ``Low-rank Matrix Recovery from Noisy, Quantized and Erroneous Measurements.'' \emph{IEEE Transactions on Signal Processing}, 2018, 66 (11): 2918-2932.}
	
\item \href{}{{\bf P. Gao}, M. Wang, J. H. Chow, M. Berger, and L. M. Seversky. ``Missing Data Recovery for High-dimensional Signals with Nonlinear Low-dimensional Structures.'' \emph{IEEE Transactions on Signal Processing}, 2017, 65 (20): 5421-5436.}
	
\item \href{}{{\bf P. Gao}, M. Wang, J. H. Chow, S. G. Ghiocel, B. Fardanesh, G. Stefopoulos, and M. P. Razanousky. ``Identification of Successive ``Unobservable'' Cyber Data Attacks in Power Systems Through Matrix Decomposition.'' \emph{IEEE Transactions on Signal Processing}, 2016, 64 (21): 5557-5570.}
		
\item \href{}{{\bf P. Gao}, M. Wang, S. G. Ghiocel, J. H. Chow, B. Fardanesh, and G. Stefopoulos. ``Missing Data Recovery by Exploiting Low-dimensionality in Power System Synchrophasor \\ Measurements.'' \emph{IEEE Transactions on Power Systems}, 2016, 31 (2): 1006-1013.}
	
\end{bibsection}


\section{Conference Publications}
\vspace{-.1275in}
\begin{bibsection}

\item \href{}{H. Qu$^*$, Y. Liu$^*$, R. Jin, W. Zhang, {\bf P. Gao}, W. Liu, and J. Luan. ``Scaling, Benchmarking, and Reasoning of Vision-Language Agents for Mobile GUI Navigation.'' \emph{Proc. of the 43rd International Conference on Machine Learning (ICML)}, 2026.}

\item \href{}{Y. Liu, W. Xu, K. Huang, C. Chen, J. Zhao, {\bf P. Gao}, W. Liu, J. Luan, S. Shang, B. Du, J. Wen, and R. Yan. ``CoME: Empowering Channel-of-Mobile-Experts with Informative Hybrid-Capabilities Reasoning.'' \emph{Proc. of the 43rd International Conference on Machine Learning (ICML)}, 2026.}

\item \href{}{Q. Wu$^*$, Z. Yang$^*$, H. Li, {\bf P. Gao}, W. Liu, and J. Luan. ``MobileBench-OL: A Comprehensive Chinese Benchmark for Evaluating Mobile GUI Agents in Real-World Environment.'' \emph{Findings of the 64th Annual Meeting of the Association for Computational Linguistics (ACL)}, 2026.}

\item \href{}{Y. Chen, Y. Liu, L. Zhang, {\bf P. Gao}, J. Luan, and W. Liu. ``STEP: Success-Rate-Aware Trajectory-Efficient Policy Optimization.'' \emph{Findings of the 64th Annual Meeting of the Association for Computational Linguistics (ACL)}, 2026.}

\item \href{}{R. Jin, {\bf P. Gao}, Y. Ren, Z. Han, T. Zhang, W. Huang, W. Liu, J. Luan, and D. Xiong. ``Revisiting Entropy in Reinforcement Learning for Large Reasoning Models.'' \emph{Findings of the 64th Annual Meeting of the Association for Computational Linguistics (ACL)}, 2026.}

\item \href{}{G. Liu, J. Ye, J. Liu, W. Liu, {\bf P. Gao}, J. Luan, Y. Li, and Y. Liu. ``Mobile GUI Agents under Real-world Threats: Are We There Yet?'' \emph{Proc. of the 24th ACM International Conference on Mobile Systems, Applications, and Services (MobiSys)}, 2026.}

\item \href{}{K. Huang, W. Xu, Y. Liu, Q. Wang, {\bf P. Gao}, W. Liu, J. Luan, B. Wang, and B. An. ``MobileIPL: Enhancing Mobile Agents Thinking Process via Iterative Preference Learning.'' \emph{Proc. of the 14th International Conference on Learning Representations (ICLR)}, 2026.}

\item \href{}{W. Xu, Z. Jiang, Y. Liu, {\bf P. Gao}, W. Liu, J. Luan, Y. Li, Y. Liu, B. Wang, and B. An. ``SMAN-Bench: A Cross-System Benchmark for Mobile Agents under Single- and Multi-path, Ambiguous, and Noisy Tasks.'' \emph{Proc. of the 14th International Conference on Learning Representations (ICLR)}, 2026.}

\item \href{}{L. Gao, L. Zhang, {\bf P. Gao}, W. Liu, J. Luan, and M. Xu. ``GUI-Shift: Enhancing VLM-Based GUI Agents through Self-supervised Reinforcement Learning.'' \emph{Proc. of the 14th International Conference on Learning Representations (ICLR)}, 2026.}

\item \href{}{G. Liu, J. Ye, J. Liu, Y. Li, W. Liu, {\bf P. Gao}, J. Luan, and Y. Liu. ``Hijacking JARVIS: Benchmarking Mobile GUI Agents against Unprivileged Third Parties.'' \emph{Proc. of the Workshop on Edge and Mobile Foundation Models (EdgeFM)}, 2025.}

\item \href{}{Q. Wu, {\bf P. Gao}, W. Liu, and J. Luan. ``BacktrackAgent: Enhancing GUI Agent with Error Detection and Backtracking Mechanism.'' \emph{Proc. of the 2025 Conference on Empirical Methods in Natural Language Processing (EMNLP)}, 2025.}

\item \href{}{M. Cui$^*$, {\bf P. Gao}$^*$, W. Liu, J. Luan, and B. Wang. ``Multilingual Machine Translation with Open Large Language Models at Practical Scale: An Empirical Study.'' \emph{Proc. of the 2025 Annual Conference of the Nations of the Americas Chapter of the Association for Computational Linguistics (NAACL)}, 2025.}

\item \href{}{M. Sun, W. Liu, J. Luan, {\bf P. Gao}, and B. Wang. ``Mixture of Diverse Size Experts.'' \emph{Proc. of the 2024 Conference on Empirical Methods in Natural Language Processing (EMNLP): Industry Track}, 2024.}

\item \href{}{J. Du, Y. Wang, W. Zhao, Z. Deng, S. Liu, R. Lou, H. P. Zou, P. N. Venkit, N. Zhang, M. Srinath, H. R. Zhang, V. Gupta, Y. Li, T. Li, F. Wang, Q. Liu, T. Liu, {\bf P. Gao}, C. Xia, C. Xing, J. Cheng, Z. Wang, Y. Su, R. S. Shah, R. Guo, J. Gu, H. Li, K. Wei, Z. Wang, L. Cheng, S. Ranathunga, M. Fang, J. Fu, F. Liu, R. Huang, E. Blanco, Y. Cao, R. Zhang, P. S. Yu, and W. Yin. ``LLMs assist NLP Researchers: Critique Paper (Meta-)Reviewing.'' \emph{Proc. of the 2024 Conference on Empirical Methods in Natural Language Processing (EMNLP)}, 2024.}

\item \href{}{{\bf P. Gao}, R. Zhang, Z. He, H. Wu, and H. Wang. ``An Empirical Study of Consistency Regularization for End-to-End Speech-to-Text Translation.'' \emph{Proc. of the 2024 Conference of the North American Chapter of the Association for Computational Linguistics (NAACL)}, 2024.}

\item \href{}{{\bf P. Gao}, L. Zhang, Z. He, H. Wu, and H. Wang. ``Learning Multilingual Sentence Representations with Cross-lingual Consistency Regularization.'' \emph{Proc. of the 2023 Conference on Empirical Methods in Natural Language Processing (EMNLP): Industry Track}, 2023.}

\item \href{}{{\bf P. Gao}, L. Zhang, Z. He, H. Wu, and H. Wang. ``Improving Zero-shot Multilingual Neural Machine Translation by Leveraging Cross-lingual Consistency Regularization.'' \emph{Findings of the 61st Annual Meeting of the Association for Computational Linguistics (ACL)}, 2023.}

\item \href{}{{\bf P. Gao}, Z. He, H. Wu, and H. Wang. ``Bi-SimCut: A Simple Strategy for Boosting Neural Machine Translation.'' \emph{Proc. of the 2022 Conference of the North American Chapter of the Association for Computational Linguistics (NAACL)}, 2022.}

\item \href{}{J. Li, {\bf P. Gao}, X. Wu, Y. Feng, Z. He, H. Wu, and H. Wang. ``Mixup Decoding for Diverse Machine Translation.'' \emph{Findings of the 2021 Conference on Empirical Methods in Natural Language Processing (EMNLP)}, 2021.}

\item \href{}{R. Wang, T. Chen, Z. Xu, and {\bf P. Gao}. ``Robust Low-Rank Tensor Recovery From Quantized and Corrupted Measurements.'' \emph{Proc. of Asilomar Conference on Signals, Systems, and Computers}, 2021.}
	
\item \href{}{Z. Liu, G. Ding, A. Bukkittu, M. Gupta, {\bf P. Gao}, A. Ahmed, S. Zhang, X. Gao, S. Singhavi, L. Li, W. Wei, Z. Hu, H. Shi, X. Liang, T. Mitamura, E. P. Xing, and Z. Hu. ``A Data-Centric Framework for Composable NLP Workflows.'' \emph{Proc. of the 2020 Conference on Empirical Methods in Natural Language Processing (EMNLP): System Demonstrations}, 2020.}
	
\item \href{}{M. Wang, J. H. Chow, Y. Hao, S. Zhang, W. Li, R. Wang, {\bf P. Gao}, C. Lackner, E. Farantatos, and M. Patel. ``A Low-rank Framework of PMU Data Recovery and Event Identification.'' \emph{Proc. of the First IEEE International Conference on Smart Grid Synchronized Measurements and Analytics (SGSMA)}, 2019.}

\item \href{}{G. M. De Mijolla, S. Konstantinopoulos, {\bf P. Gao}, J. H. Chow, and M. Wang. ``An Evaluation of Algorithms for Synchrophasor Missing Data Recovery.'' \emph{Proc. of the 20th Power Systems Computation Conference (PSCC)}, 2018.}

\item \href{}{{\bf P. Gao} and M. Wang. ``Dynamic Matrix Recovery from Partially Observed and Erroneous Measurements.'' \emph{Proc. of the International Conference on Acoustics, Speech and Signal Processing (ICASSP)}, 2018.}
	
\item \href{}{M. Wang, J. H. Chow, {\bf P. Gao}, Y. Hao, W. Li, and R. Wang. ``Recent Results of PMU Data Analytics by Exploiting Low-dimensional Structures.'' \emph{Proc. of the 10th Bulk Power Systems Dynamics and Control Symposium (IREP)}, 2017.}
	
\item \href{}{{\bf P. Gao}, R. Wang, and M. Wang. ``Quantized Low-rank Matrix Recovery with Erroneous Measurements: Application to Data Privacy in Power Grids.'' \emph{Proc. of Asilomar Conference on Signals, Systems, and Computers}, 2016.}
	
\item \href{}{{\bf P. Gao}, M. Wang, and J. H. Chow. ``Matrix Completion with Columns in Union and Sums of Subspaces.'' \emph{Proc. of IEEE Global Conference on Signal and Information Processing (GlobalSIP)}, 2015.}
	
\item \href{}{M. Wang, J. H. Chow, {\bf P. Gao}, X. T. Jiang, Y. Xia, S. G. Ghiocel, B. Fardanesh, G. Stefopoulos, Y. Kokai, N. Saito, and M. P. Razanousky. ``A Low-Rank Matrix Approach for the Analysis of Large Amounts of Power System Synchrophasor Data.'' \emph{Proc. of Hawaii International Conference on System Sciences (\textcolor{red}{Runner-up of Best Paper in Electric Energy Systems Track})}, 2015.}
	
\item \href{}{M. Wang, {\bf P. Gao}, S. G. Ghiocel, J. H. Chow, B. Fardanesh, G. Stefopoulos, and M. P. Razanousky. ``Identification of ``Unobservable'' Cyber Data Attacks on Power Grids.'' \emph{Proc. of IEEE SmartGridComm}, 2014.}
	
\item \href{}{{\bf P. Gao}, M. Wang, S. G. Ghiocel, and J. H. Chow. ``Modeless Reconstruction of Missing Synchrophasor Measurements.'' \emph{Proc. of IEEE Power \& Energy Society General Meeting (\textcolor{red}{Best Conference Papers on Power System Analysis and Modeling})}, 2014.}

\end{bibsection}


% \section{Technical Reports}
% \vspace{-.1275in}
% \begin{bibsection}
	
% \item \href{}{Z. Hu, {\bf P. Gao}, A. Bukkittu, and Z. Hu. ``Introducing Texar-PyTorch: An ML Library integrating the best of TensorFlow into PyTorch." October, 2019.}

% \end{bibsection}


% \section{Patents}
% \vspace{-.1275in}
% \begin{bibsection}

% \item \href{}{{\bf P. Gao}, R. Zhang, Z. He, and H. Wu. ``Method and device for training speech translation model, and storage medium.'' Application No.: US18/930,081, Filed October 29, 2024.}

% \item \href{}{{\bf P. Gao}, Z. He, and H. Wu. ``\begin{CJK}{UTF8}{gbsn}用于多语言翻译的语言模型的训练方法及翻译方法.\end{CJK}'' Application No.: CN202311759677.9, Filed December 20, 2023.}

% \item \href{}{{\bf P. Gao}, R. Zhang, Z. He, and H. Wu. ``\begin{CJK}{UTF8}{gbsn}语音翻译模型的训练方法、语音翻译方法和装置.\end{CJK}'' Application No.: CN202311622229.4, Filed November 30, 2023.}

% \item \href{}{{\bf P. Gao}, L. Zhang, Z. He, Z. Li, and H. Wu. ``\begin{CJK}{UTF8}{gbsn}多语言的句向量的获取方法和多语言的编码器的训练方法.\end{CJK}'' Application No.: CN202310395274.4, Filed April 13, 2023.}

% \item \href{}{{\bf P. Gao}, Z. He, Z. Li, and H. Wu. ``\begin{CJK}{UTF8}{min}深層学習モデルのトレーニング方法及び装置、テキストデータ処理方法及び装置、電子機器、記憶媒体、並びにコンピュータプログラム.\end{CJK}'' Application No.: JP2022-190230, Filed November 29, 2022.}

% \item \href{}{{\bf P. Gao}, Z. He, Z. Li, and H. Wu. ``Method of training deep learning model and method of processing text data.'' Application No.: US18/059,389, Filed November 28, 2022.}

% \item \href{}{{\bf P. Gao}, Z. He, Z. Li, and H. Wu. ``\begin{CJK}{UTF8}{gbsn}深度学习模型的训练方法、文本数据处理方法和装置.\end{CJK}'' Application No.: CN202210189268.9, Filed February 28, 2022.}

% \item \href{}{{\bf P. Gao}, Z. He, Z. Li, and H. Wu. ``\begin{CJK}{UTF8}{gbsn}一种模型训练方法、装置、电子设备及存储介质.\end{CJK}'' Application No.: CN202210186688.1, Filed February 28, 2022.}

% \item \href{}{{\bf P. Gao}, Z. He, H. Wu, and H. Wang. ``\begin{CJK}{UTF8}{gbsn}用于训练模型的方法、装置、设备、介质和程序产品.\end{CJK}'' Application No.: CN202111288550.4, Filed November 2, 2021.}

% \item \href{}{X. Wan, J. Zhao, M. Wang, Z. He, H. Wu, Z. Li, Z. Xu, J. Liu, {\bf P. Gao}, M. Sun, C. Li, and W. Yao. ``\begin{CJK}{UTF8}{gbsn}翻译模型的训练方法、装置、电子设备及存储介质.\end{CJK}'' Application No.: CN202111014476.7, Filed August 31, 2021.}

% \item \href{}{J. Li, {\bf P. Gao}, Z. He, and Z. Li. ``\begin{CJK}{UTF8}{gbsn}文本翻译方法、装置、电子设备及存储介质.\end{CJK}'' Application No.: CN202110736794.8, Filed June 30, 2021.}

% \item \href{}{J. Li, {\bf P. Gao}, Z. He, and Z. Li. ``\begin{CJK}{UTF8}{gbsn}语料生成方法、装置、电子设备以及存储介质.\end{CJK}'' Application No.: CN202110748376.0, Filed June 30, 2021.}

% \item \href{}{{\bf P. Gao}, Z. He, H. Wu, and H. Wang. ``\begin{CJK}{UTF8}{gbsn}模型训练方法、装置、电子设备及计算机可读存储介质.\end{CJK}'' Application No.: CN202110320138.X, Filed March 25, 2021.}

% \item \href{}{M. Wang, {\bf P. Gao}, and J. H. Chow. ``A low-rank-based missing PMU data recovery method.'' Application No.: 62/445305, Filed January 12, 2017.}
	
% \end{bibsection}


%\section{Research Experience}

%\textbf{Research Assistant} \hfill {August 2013 to December 2017}
%\begin{innerlist}
%\item[] ECSE Department,\\
%\textbf{Rensselaer Polytechnic Institute}\\
%Supervisor: Professor Meng Wang
%\item Analyzed the Phasor Measurement Unit (PMU) data ($>$ 200 MB of data) to exploit the temporal and spatial correlations (low dimensionality) of the data.
%\item Proposed an identification method that can detect the cyber data attack in the power system. Tested our method on the actual PMU data from Central New York Power System.
%\item Developed an on-line algorithm to estimate the missing PMU data in real time manner. Built the corresponding action adapter in OpenPDC by C\# code, reducing the computational time by 50\%.
%\item Proposed a novel model to characterize the practical nonlinear datasets. Developed convex-optimization-based methods to recover missing data under this model. Tested our method on simulated power system data in IEEE 39-bus New England Power System. 
%\item Proposed a novel method to recover the original data from quantized measurements even when partial measurements are corrupted. Developed a projected gradient method to solve the non-convex problem approximately. Tested our method on actual PMU data from Central New York Power System. \\
%\end{innerlist}

%\textbf{Member of RPI Team for IBM Watson Build Challenge} \hfill {July 2017 to October 2017}
%\begin{innerlist}
%\item[] ECSE Department,\\
%\textbf{Rensselaer Polytechnic Institute}\\
%Supervisor: Professor Meng Wang
%\item Developed a web application with IBM Watson APIs based on IBM Bluemix. 
%\item Implemented multiple missing data recovery algorithms by Python for large scale Phasor Measurement Unit (PMU) data analysis.\\
%\end{innerlist}

%\textbf{Summer Research Co-program Manager} \hfill {May 2016 to August 2016}
%\begin{innerlist}
%\item[] ECSE Department,\\
%\textbf{Rensselaer Polytechnic Institute}\\
%Co-program Manager: Stephen M. Burchett
%\item Led a group of 3 undergraduates to build a wireless and real-time photovoltaic power monitoring system. The data acquisition was performed by the Arduino platform. Data communication was achieved by Bluetooth connection. And the monitoring interface was developed by Matlab.\\
%\end{innerlist}

%\textbf{Research Assistant} \hfill {May 2012 to May 2013}
%\begin{innerlist}
%\item[] Department of Bioengineering,\\
%\textbf{University of Pennsylvania}\\
%Supervisor: Professor Gershon Buchsbaum        
%\item Analyzed the EEG data from IEEG Portal for epilepsy detection.
%\item Proposed a new dictionary for EEG dataset. Improved the reconstruction performance of the EEG data by 20\%. \\
%\end{innerlist}

\section{Projects}

\textbf{Forte: A Data-Centric Framework for Composable NLP Workflows} 
\begin{innerlist}	
\item[] Core Machine Learning Team,\\
\textbf{Petuum, Inc.}, Pittsburgh, PA, USA
\item Forte is a toolkit for building NLP pipelines, featuring composable components, convenient data interfaces, and cross-task interaction. Forte designs a universal data representation format for text, making it a one-stop platform to assemble state-of-the-art NLP/ML technologies, ranging from Information Retrieval and Natural Language Understanding to Natural Language Generation. Forte was originally developed at CMU and is actively contributed by Petuum in collaboration with other institutes. I am one of the main contributors to Forte.
\end{innerlist}

\textbf{Texar-PyTorch: A Modularized, Versatile, and Extensible Toolkit for Text Generation} 
\begin{innerlist}	
\item[] Core Machine Learning Team,\\
\textbf{Petuum, Inc.}, Pittsburgh, PA, USA
\item Texar-PyTorch is a toolkit aiming to support a broad set of machine learning, especially natural language processing and text generation tasks. Texar-PyTorch provides a library of easy-to-use ML modules and functionalities for composing whatever models and algorithms. The tool is designed for both researchers and practitioners for fast prototyping and experimentation. Texar-PyTorch was originally developed and is actively contributed by Petuum and CMU in collaboration with other institutes. I am one of the main contributors to Texar-PyTorch.
\end{innerlist}

%\textbf{DyNet: The Dynamic Neural Network Toolkit} 
%\begin{innerlist}	
%\item[] Machine Learning Team,\\
%\textbf{Petuum, Inc.}, Pittsburgh, PA, USA
%\item DyNet is a neural network library developed by CMU, Petuum, and many others. It is written in C++ (with bindings in Python) and is designed to be efficient when run on either CPU or GPU, and to work well with networks that have dynamic structures that change for every training instance. I contribute the example and tutorial part of DyNet repository.\\
%\end{innerlist}

%\textbf{IBM Watson Build Challenge 2017}
%\begin{innerlist}
%\item[] ECSE Department,\\
%\textbf{Rensselaer Polytechnic Institute}
%\item We developed a web application with IBM Watson APIs based on IBM Bluemix. We implemented multiple missing data recovery algorithms in Python for large-scale Phasor Measurement Unit (PMU) data analysis.\\
%\end{innerlist}

%\textbf{Online Algorithm for PMU Data Processing (OLAP)} 
%\begin{innerlist}	
%\item[] ECSE Department,\\
%\textbf{Rensselaer Polytechnic Institute}	
%\item We implemented OLAP by C\# based on Project Alpha for the real-time application. Project Alpha is the elite version of Open PDC. The code developed on Project Alpha can be run on Open PDC as an action adapter.\\
%\end{innerlist}

%\textbf{Mobile Eye Gaze Estimation with Deep Learning} 
%\begin{innerlist}	
%\item[] ECSE Department,\\
%\textbf{Rensselaer Polytechnic Institute}	
%\item We implemented a deep convolutional neural network based on TensorFlow for eye gaze estimation. The original dataset comes from the GazeCapture project. Our model follows the architecture introduced in CVPR paper ``Eye Tracking for Everyone'', and we changed and tuned the hyper parameters due to the different image size in our training dataset.\\
%\end{innerlist}

%\textbf{Twitter Sentiment Analysis with Recurrent Neural Networks} 
%\begin{innerlist}	
%\item[] ECSE Department,\\
%\textbf{Rensselaer Polytechnic Institute}	
%\item We implemented a recurrent neural network based on TensorFlow for the task of sentiment analysis on natural language data. We analyzed the data from Twitter (Sentiment140 dataset) and classified it as either “positive” or “negative”. 
%\end{innerlist}

\section{Professional Activities \& Service }
\begin{innerlist}
\item Member of IEEE, 2018 - present. Member of ACL, 2022 - present.
% \item RPI Student Representative at the Center for Ultra-wide-area Resilient Electric Energy Transmission Networks (CURENT), 2015 - 2016. 
% \item Teaching Assistant (Rensselaer Polytechnic Institute): \\
% Modeling and Analysis of Uncertainty, Fall 2017, \\
% Distributed Systems and Sensor Networks, Fall 2017.
\item Area Chair: \\
The Association for Computational Linguistics Rolling Review (ACL ARR).
\item Journal Reviewer: \\
IEEE Transactions on Signal Processing, \\
IEEE Transactions on Audio, Speech and Language Processing, \\
IEEE Transactions on Smart Grid, \\
IEEE Transactions on Automatic Control, \\
IEEE Transactions on Power Delivery, \\
IEEE/ACM Transactions on Networking, \\
IEEE Signal Processing Letters, \\
Annals of Mathematics and Artificial Intelligence, \\
Journal of Data-centric Machine Learning Research (DMLR).
\item Program Committee Member / Conference Reviewer: \\
The Conference on Uncertainty in Artificial Intelligence (UAI), \\
The Annual AAAI Conference on Artificial Intelligence (AAAI), \\
The Conference on Neural Information Processing Systems (NeurIPS), \\
The International Conference on Learning Representations (ICLR), \\
The International Conference on Machine Learning (ICML), \\
The Association for Computational Linguistics Rolling Review (ACL ARR), \\
The Annual Meeting of the Association for Computational Linguistics (ACL), \\
The Conference on Empirical Methods in Natural Language Processing (EMNLP), \\
The Annual Conference of the North American Chapter of the Association for Computational Linguistics (NAACL), \\
The Conference of the European Chapter of the Association for Computational Linguistics (EACL), \\
The International Conference on Computational Linguistics (COLING), \\
The Conference on Language Modeling (COLM), \\
The China National Conference on Computational Linguistics (CCL), \\
The CCF International Conference on Natural Language Processing and Chinese Computing (NLPCC), \\
SIAM International Conference on Data Mining (SDM), \\
The ACM SIGKDD Conference on Knowledge Discovery and Data Mining (KDD), \\
IEEE/CVF Conference on Computer Vision and Pattern Recognition (CVPR), \\
IEEE/CVF International Conference on Computer Vision (ICCV), \\
The European Conference on Computer Vision (ECCV), \\
British Machine Vision Conference (BMVC), \\
IEEE International Conference on Communications, Control, and Computing Technologies for Smart Grids (SmartGridComm), \\
American Control Conference (ACC), \\
Intelligent System Applications to Power Systems Conference (ISAP), \\
International Conference on Intelligent Multilingual Information Processing (IMLIP).


\end{innerlist}

%\halfblankline
%
%\section{Posters \& Presentations}
%\begin{innerlist}
%\item Asilomar Conference on Signals, Systems, and Computers \hfill Nov 2016
%\item IEEE Global Conference on Signal and Information Processing \hfill Dec 2015
%\item CURENT's 2015 Industry Conference \& NSF/DOE Site Visit \hfill Oct 2015
%\item CFES Annual Conference \hfill Feb 2015
%\item IEEE Power \& Energy Society General Meeting \hfill July 2014
%\item CURENT's 2014 Industry Conference \& NSF/DOE Site Visit \hfill May 2014
%\item CFES Annual Conference \hfill Feb 2014
%\end{innerlist}

%\halfblankline

\end{document}

%%%%%%%%%%%%%%%%%%%%%%%%%% End CV Document %%%%%%%%%%%%%%%%%%%%%%%%%%%%%

%----------------------------------------------------------------------%
% The following is copyright and licensing information for
% redistribution of this LaTeX source code; it also includes a liability
% statement. If this source code is not being redistributed to others,
% it may be omitted. It has no effect on the function of the above code.
%----------------------------------------------------------------------%
% Copyright (c) 2007, 2008, 2009, 2010, 2011 by Theodore P. Pavlic
%
% Unless otherwise expressly stated, this work is licensed under the
% Creative Commons Attribution-Noncommercial 3.0 United States License. To
% view a copy of this license, visit
% http://creativecommons.org/licenses/by-nc/3.0/us/ or send a letter to
% Creative Commons, 171 Second Street, Suite 300, San Francisco,
% California, 94105, USA.
%
% THE SOFTWARE IS PROVIDED "AS IS", WITHOUT WARRANTY OF ANY KIND, EXPRESS
% OR IMPLIED, INCLUDING BUT NOT LIMITED TO THE WARRANTIES OF
% MERCHANTABILITY, FITNESS FOR A PARTICULAR PURPOSE AND NONINFRINGEMENT.
% IN NO EVENT SHALL THE AUTHORS OR COPYRIGHT HOLDERS BE LIABLE FOR ANY
% CLAIM, DAMAGES OR OTHER LIABILITY, WHETHER IN AN ACTION OF CONTRACT,
% TORT OR OTHERWISE, ARISING FROM, OUT OF OR IN CONNECTION WITH THE
% SOFTWARE OR THE USE OR OTHER DEALINGS IN THE SOFTWARE.
%----------------------------------------------------------------------%
